\section{Introducción}

% Introducción a nanoportículas y nanotecnología
En las últimas décadas, el interés por la nanotecnología ha incrementado constantemente, impulsado por la posibilidad de explotar las propiedades físicas y  químicas que solo se presentan a escala nanométrica. % Introduce el tema
En este contexto, las nanopartículas (NPs) han ocupado un papel central en el campo: al poseer alta relación volumen-superficie emergen propiedades marcadamente distintas a las de los mismos materiales en estado bulk. % Interés en las nanopartículas (No falta decir que son las nanopartículas?)
Estudiar estas estructuras desde primeros principios permite determinar los principales factores tanto en su formación como en sus propiedades funcionales.

% Nanopartículas de plata y oro en general
Entre las NPs que han recibido mayor atención están las de oro y plata, dado que permiten crear sistemas con propiedades ópticas y electrónicas singulares, con aplicaciones en biosensado, catálisis y fototerapia, entre otras. % Mejorar la conexión entre esta oración y la siguiente.
A escala nano, las propiedades son altamente dependientes la geometría de las partículas.
Por esto, gran parte de la investagación en el campo se oriente a desarrollar mecanismos de síntesis que permitan controlar su forma con precisión.

% Nanopartículas triangulares de plata (Importante!)
Este trabajo se centra en los nanoplatos triangulares de plata (AgTNPs), estructuras cuyo ancho es ampliamente mayor a su espesor.
Su geometría particular las hace especialmente sensibles al campo electromagnético: presentan resonancia plasmónica superficial (SPR), un fenómeno en el que los electrones de conducción oscilan colectivamente en respuesta a la radiación incidente, con una frecuencia sintonizable mediante el tamaño y la forma de la partícula \cite{}. 
Tales propiedades hacen a las nanopartículas triangulares posibles candidatos para medir potencial químico y campo eléctrico de manera localizada.
Además, en los vértices el campo electromagnético se ve amplificado por efecto punta.
De manera que, al ser excitadas con un radiación externa, las partículas cercanas a las esquinas reciben una gran patadón de energia.
Haciendo a las AgTNPs candidatos para el desarrollo de catalizadores.
Finalmente, la plata es relativamente económica, por lo que utilización de manera industrial no ha sido descartada.

Controlar las dimensiones de las AgTNPs permite sintonizar su frecuencia de resonancia plasmónica con gran precisión, lo que hace indispensable entender los procesos involucrados en su crecimiento.
En relación a esto, se ha visto que la formación de nanopartículas triangulares está altamente vinculada a la presencia de fallas de apilamiento en el eje {111} \cite{tan2021}.
Además, se ha observado que aunque el ancho de las nanopartículas suele variar, el grosor se mantiene uniforme.
¿Qué causa que el grosor se mantenga constante? y ¿cuál es el rol de las fallas en el crecimiento de las nanopartículas? son preguntas que todavía están en constante discusión. 
Para avanzar en su estudio sería muy útil caracterizar la cantidad, presencia y disposición de fallas a lo largo de una nanopartícula. 

Aquí La microscopia de transmisión electrónica (TEM) permite atacar el problema.
Mediante TEM se estudia de manera localizada el ordenamiento de los diferentes planos y átomos dentro de la partícula.
Si se observa una partícula de canto, es trivial encontrar la cantidad y tipo de fallas de apilamiento. 
Sin embargo, al suspender las nanopartículas sobre la grilla del portamuestra, la mayoría se acomoda con las caras planas {111} sobre el eje óptico.
Dada la escasa cantidad de nanopartículas que se acomodan sobre el canto, es complicado obtener un resultado representativo de toda la muestra simplemente estudiando las partículas que encontramos à côté.

En este trabajo, se propone un mecanismo innovador para estudiar la disposición de las fallas en la NPs simplemente estudiando la cara {111} de las NPs.
Se realizan simulaciones multislice de imágenes TEM para un conjunto de cantidad y posición de fallas dentro de la pila de planos.
A partir las simulaciones, se compara con las imágenes experimentales para obtener información de los posibles ordenamientos de los átomos asociados a la nanopartícula fotografiada.
De esta manera, se logra realizar un estudio amplio sobre la disposición de las fallas dentro las NPs.